\chapter{Introduction}

\section{Background}

Presentations are an essential part of how we communicate ideas in colleges, universities, corporate offices, and technical conferences. For decades, computerized presentation tools like Microsoft PowerPoint and Google Slides have served as the standard way to prepare visual aids. While these tools make it easy to format text, add graphics, and show bullet points, they are fundamentally designed for one-way broadcasting. The presenter speaks, the slides advance, and the audience remains passive spectators.

In modern educational and professional settings, passive listening often results in short attention spans and low retention rates. Educational research has repeatedly shown that active learning—where listeners participate through questions, polls, and discussions—greatly improves understanding and keeps participants focused.

To make sessions more engaging, many speakers try using third-party polling tools such as Mentimeter, Slido, or Kahoot. However, because these tools are separate from standard presentation software, presenters must constantly switch between their presentation slides and a web browser window during live talks. Audience members must also visit separate websites, enter room codes, and wait for popups. This constant back-and-forth disrupts the speaker's rhythm and distracts the audience.

\textbf{Sentio} was created to eliminate this friction. Sentio is a unified, web-based presentation platform that integrates visual slide authoring, live audience polling, competitive quizzes, and AI assistance into a single, cohesive tool.

\section{Problem Statement}

Presenters and audiences regularly face several practical difficulties during live presentations:

\begin{enumerate}
    \item \textbf{Passive and Disengaged Audiences:} Traditional slide tools provide no built-in way for listeners to participate, ask questions, or share feedback during a talk.
    \item \textbf{Tool Switching and Presentation Friction:} Juggling separate applications for slides, live polls, and Q\&A creates awkward pauses, screen flickering, and technical confusion during presentations.
    \item \textbf{Rigid Question Layouts:} Most polling tools display questions as fixed full-screen takeovers, preventing presenters from combining questions with explanatory text, code snippets, or diagrams on the same slide.
    \item \textbf{High Preparation Overhead:} Manually brainstorming and drafting interactive quiz questions, distractors, and discussion topics takes considerable time and effort.
\end{enumerate}

Sentio directly addresses these challenges by offering an all-in-one presentation environment where interactive elements sit directly on the slide canvas, synchronized in real time via WebSockets and supported by fast generative AI.

\section{Objectives}

The core objectives of the Sentio project are:

\begin{itemize}
    \item \textbf{Develop a 16:9 Canvas Slide Editor:} Build a flexible visual editor that maintains a reliable 16:9 aspect ratio across all screen sizes without shifting elements or clipping content.
    \item \textbf{Integrate Live Interaction Elements:} Allow presenters to place live quizzes, multiple-choice polls, rating matrices, word clouds, Q\&A boards, and leaderboards directly on any slide.
    \item \textbf{Provide Sub-Second Real-Time Synchronization:} Implement a high-speed WebSocket engine using Socket.IO so audience responses, timer updates, and slide transitions happen with near-zero latency.
    \item \textbf{Embed Fast AI Slide Generation:} Integrate Groq AI with the LLaMA 3.3 model to automatically create complete slide decks, quiz questions, and summaries from simple topic prompts or uploaded documents.
    \item \textbf{Deliver Role-Specific Views:} Create tailored interfaces for Presenters (with private speaker notes and live controls), Projectors (a clean display without clutter), and Participants (a touch-friendly mobile interface).
    \item \textbf{Automate Post-Session Analytics:} Provide detailed session attendance, response breakdowns, and downloadable PDF reports.
\end{itemize}

\section{Scope}

The scope of the Sentio platform covers:

\begin{itemize}
    \item \textbf{Slide Studio:} An intuitive 3-panel builder with slide thumbnail navigation, a visual drag-and-drop canvas, and a context-sensitive properties panel with automatic saving.
    \item \textbf{Audience Participation:} A lightweight mobile web client that allows audience members to join instantly using a 6-character PIN or QR code without requiring an account.
    \item \textbf{Live Host Controls:} Presenter tools for starting timers, locking responses, revealing correct answers, and displaying real-time leaderboards.
    \item \textbf{User and Organization Management:} Secure account registration, JWT authentication, password encryption, and multi-user organization workspaces.
    \item \textbf{Reporting:} Instant generation of clean, downloadable PDF summary reports at the conclusion of every presentation.
\end{itemize}

\section{Relevance}

Sentio brings clear value to modern presentations:

\begin{itemize}
    \item \textbf{For Educators and Trainers:} Helps instructors gauge student understanding on the spot and immediately address topics where students struggle.
    \item \textbf{For Speakers and Presenters:} Eliminates the need to buy and manage multiple software subscriptions by unifying slides and polling in one application.
    \item \textbf{From a Software Engineering Standpoint:} Combines modern full-stack web technologies—including Next.js 15, React 19, TypeScript, Socket.IO, and Groq AI—into a maintainable monorepo architecture.
\end{itemize}

\section{Report Organization}

The report is structured into the following chapters:

\begin{itemize}
    \item \textbf{Chapter 2 (System Study):} Examines existing systems, analyzes their limitations, introduces the proposed platform, and details the feasibility study.
    \item \textbf{Chapter 3 (Development Methodology):} Explains the Agile Scrum methodology, sprint planning, and the justification for the chosen technology stack.
    \item \textbf{Chapter 4 (System Design):} Presents the architectural blueprints, data flow diagrams, UML diagrams, and database schemas.
    \item \textbf{Chapter 5 (Implementation and Results):} Covers the development of the frontend and backend modules, testing procedures, and system results.
    \item \textbf{Chapter 6 (Conclusion and Future Scope):} Summarizes the project achievements, notes current limitations, and outlines future enhancements.
\end{itemize}
